% =====================================================================
%  ABSTRACT
% =====================================================================
\chapter*{Abstract}
\addcontentsline{toc}{chapter}{Abstract}

\noindent At the Universidad Nacional de San Antonio Abad del Cusco (UNSAAC), almost no experience of sexual harassment ends up as a formal complaint. Many of these situations are recorded only as written narratives, and recognizing harassment in them requires careful reading, since it is not always described in explicit terms. This research proposes to develop a machine learning model that, given a narrative written by a student, determines whether it describes a situation of sexual harassment. A corpus of narratives voluntarily contributed by adult students will be built, anonymized, and labeled by two annotators following a guide reviewed by specialists. Development will follow the CRISP-DM methodology. Logistic regression, support vector machine, and Naïve Bayes with TF-IDF representation will be trained as candidate models, along with the BETO model fine-tuned for Spanish. All processing will run on local machines, making sure that test data never take part in training. The chosen model will be compared with a keyword-based lexical baseline using macro F1, precision, and recall, together with McNemar's test. The research is applied and technological, with a quantitative approach and an experimental design. The outcome will be a validated model and a prototype that shows the prediction and the words that most influenced it, intended to support the staff who review these cases rather than to replace their judgment.

\vspace{0.5em}
\noindent\textit{Keywords}: sexual harassment, machine learning, natural language processing, text classification, higher education.
